problem:  
**Conjecture (Near‑Maximal Isometry‑Dimension Rigidity for Positively Curved Manifolds).**  
Let \( M^n \) be a compact, simply connected Riemannian manifold with strictly positive sectional curvature.  
Classically, the standard sphere \(S^n\) achieves the maximal possible isometry dimension \( \dim\,\mathrm{Iso}(S^n) = \frac{1}{2}n(n+1) \).  
Let us define the *symmetry deficit*  
\[
\delta(M) = \tfrac{1}{2}n(n+1) - \dim\,\mathrm{Iso}_0(M).
\]  
Known results (Kobayashi; Grove–Searle) determine the diffeomorphism type of \(M\) when \(\delta(M)=0\) (maximal) and when the **rank** of a torus subgroup in \(\mathrm{Iso}_0(M)\) is maximal.

**Problem (precisely stated):**  
Does there exist a universal constant \(C>0\), independent of \(n\), such that if  
\[
\delta(M) \le Cn,
\]
then \(M\) must be diffeomorphic to a compact rank‑one symmetric space (CROSS)?  
Equivalently, does positive sectional curvature admit only CROSS examples among manifolds whose isometry dimension is within \(O(n)\) of the maximum possible value \( \tfrac{1}{2}n(n+1) \)?  

---

Why is it a "good" problem:  
This formulation refines the symmetry‑rigidity question into a sharply quantitative conjecture that bridges **Kobayashi's upper bound** on isometry dimension and the **Grove–Searle symmetry rank classification**. It asks whether a *linear‑scale deficit in symmetry dimension* (measured by \(O(n)\) rather than \(O(n^2)\)) still forces the manifold's topology to be that of a CROSS.  

Conceptually, the conjecture probes the “gap phenomenon” near maximal symmetry in positive curvature, analogous to spectral gap problems in analysis: is the set of possible symmetry dimensions discrete near its upper edge?  
A positive resolution would give a universal gap rigidity law linking **Lie group representation geometry** (the size of the isometry group) to **global curvature topology**, unifying transformation group theory and Riemannian geometry.  
A counterexample would demand a new mechanism for building highly symmetric positively curved manifolds distinct from known CROSSes—an unexpected and transformative breakthrough.

The statement is compact, natural, and precisely one or two steps beyond established results, representing a realistic yet profound “next frontier” in the Grove symmetry program—hence, it qualifies as a *good* research‑level problem.